
\begin{obereflection}
\textbf{Latihan 10.1}
1. Sebutkan tiga penemu algoritma RSA, lembaganya, dan tahun pengenalannya.
2. Jelaskan mengapa keamanan RSA terletak pada sulitnya memfaktorkan $n$
   menjadi $p$ dan $q$.
3. Dengan $p = 47$ dan $q = 71$, hitunglah $\varphi(n)$. Periksa bahwa
   $e = 17$ memenuhi syarat sebagai kunci publik, lalu tentukan $d$ yang
   bersesuaian.
4. Jelaskan mengapa $p$, $q$, dan $\varphi(n)$ harus dirahasiakan, sedangkan
   $n$ boleh diketahui umum.
5. Dengan kunci publik $(79, 3337)$ dan kunci privat $(1019, 3337)$,
   enkripsikan $m = 16$, lalu dekripsikan hasilnya.
6. Mengapa pesan yang lebih panjang daripada modulus harus dipecah menjadi
   blok dengan syarat $0 \le m_i \le n-1$?
\end{obereflection}

\begin{obereflection}
\textbf{Latihan 10.2}
1. Jelaskan mengapa RSA lebih lambat daripada algoritma kunci-simetri, lalu
   jelaskan cara kriptografi hibrida mengatasi hal itu.
2. Uraikan jalannya serangan \textit{Man-in-the-Middle} ketika Alice dan Bob
   bertukar kunci publik.
3. Jelaskan mengapa \textit{padding} OAEP membuat cipherteks menjadi tidak
   deterministik.
4. Bahan kuliah menyatakan bahwa modulus 256 bit dapat difaktorkan dalam
   beberapa jam dan modulus 512 bit dengan beberapa ratus komputer. Apa
   implikasi pernyataan itu bagi pemilihan panjang modulus RSA?
5. Pada tahun 1977 ketiga penemu RSA membuat sayembara di \textit{Scientific
   American}. Kapan sayembara itu terpecahkan, dan berapa lama waktu yang
   dibutuhkan?
6. Jelaskan mengapa algoritma Shor mengancam RSA, dan mengapa memperbesar
   modulus tidak menangkal ancaman itu.
\end{obereflection}
